% ====================================================================
\section{AI assistance and verification; credits}\label{sec:ai}
% ====================================================================

\subsection*{AI assistance and verification}
The mathematics in this manuscript was developed by AI agents (Claude, Anthropic) in research sessions directed by the
team. The agents proposed the strategy, proved the lemmas, and wrote the project notes from which this manuscript was
assembled. The manuscript itself was also prepared by an AI agent.

The arguments were checked by \emph{independent AI referee passes}: separate agent sessions that had not written the
material under review, re-derived each step by hand, and ran their own computer tests. In summary:
\begin{itemize}[leftmargin=1.5em]
\item one pass on Theorem~\ref{thm:LR} (with Lemma~\ref{lem:R}, the cleaning lemma, Lemmas~A, B and NF), on
  Corollary~\ref{cor:RQlist} and its relation to~\cite{SingerSudan}, and on an earlier proof of the case $n=4$. Verdict:
  correct, with editorial remarks only;
\item one pass on the kill step, the two-class count and Theorem~R$^+$. Verdict: correct after fixes F1--F3, of which
  F1 repaired a genuine gap;
\item one pass on the case $n=5$ of the class-count argument. Verdict: correct with expository gaps; items A--G of that
  report are addressed in Sections~\ref{sec:framework}--\ref{sec:core};
\item one pass on the general-$n$ argument (configurations, HR, kill step, Core Lemma, Slicing Lemma, threshold tower).
  Verdict: correct with expository gaps; items X1--X6 of that report are addressed in Sections~\ref{sec:prelim},
  \ref{sec:framework}--\ref{sec:main} and~\ref{sec:concl}.
\end{itemize}
Each referee found the material under review correct, conditional on its inputs. All of them identified
Theorem~\ref{thm:LR}, which is new and unpublished, as the weakest link.

\textbf{No human expert has reviewed any part of this argument.} AI referee passes can miss errors, and several
passes sharing similar blind spots is a real possibility. The Lean formalization of Section~\ref{sec:lean} was also
written by AI agents. It removes the dependence on the referees for the \emph{formal} statement: Lean's kernel checks
every step. Two things remain for human judgement: whether the formal statement faithfully captures Erd\H{o}s's
question, and the sharper form of Theorem~\ref{thm:LR} stated here, whose NF3 step is not formalized.

We also note what has \emph{not} been refereed in the present form. While consolidating the project notes into this
manuscript we made some routine changes, and these have not been checked by any referee pass. They are: the
explicit choices $Q=2^{44}n^4(D+2)^{8k'}$, $\smin=2^{\Lst+2+2n}(h+1)$ and $e_A=2^{\Lst+2}(h+1)$; the explicit constant
$2^{14}$ in $L_k(r)$ and the explicit cost bounds in Lemma~\ref{lem:NF}; the column-basis proof of the cleaning lemma; the
self-contained statement and proof of Lemma~\ref{lem:twodist}; the direct (rather than translated) single-cap proof of
Theorem~\ref{thm:Rplus}; and the explicit growth recursion in Lemma~\ref{lem:growth}. We believe these changes are correct
and harmless, but they are part of what an expert should check.

\subsection*{Author list and credits}
\textbf{[PLACEHOLDER --- the author list and the form of credit are to be confirmed by the team.]}

\smallskip\noindent
\emph{Direction.} Will Howe (H0W4RD-dev) directed the work.

\smallskip\noindent
\emph{Contributors and acknowledgements.} The team also includes Jensen Kohlmeyer (jenw1n), Liam Kruer (lkruer) and Grant
Huffman (Ghuffy1). They are listed here as team members; the project records available to us do not attribute any
mathematical contribution to this result to them. We acknowledge the prior in-project work of a peer session~\cite{Peer},
on which Section~\ref{sec:intro} and Lemma~\ref{lem:H} build.
